\documentclass[a4paper,12pt]{article}

\usepackage[T1]{fontenc}
\usepackage[utf8]{inputenc}
\usepackage[english,italian]{babel}
\usepackage{amsmath,amssymb}
\usepackage{newtxtext,newtxmath}          %

\usepackage[a4paper,top=3cm,bottom=3cm,left=3.5cm,right=3cm,
            head=1.35cm,headsep=0.45cm,foot=1.1cm]{geometry}

\usepackage{graphicx}
\graphicspath{{figure/}}
\usepackage{booktabs}
\usepackage{tabularx}
\usepackage{array}
\usepackage{setspace}
\usepackage{fancyhdr}
\usepackage{enumitem}
\usepackage{multicol}
\usepackage{cite}
\usepackage{etoolbox}
\usepackage[autostyle=true]{csquotes}
\usepackage[expansion=alltext,stretch=25,shrink=25,tracking=true]{microtype}
\usepackage[hidelinks]{hyperref}
\Urlmuskip=0mu plus 1mu\relax
\def\UrlBreaks{\do\/\do\-\do\.\do\_\do\?\do\=\do\&}

\patchcmd{\thebibliography}{\section*{\refname}}{\section{\refname}}{}%
  {\PackageError{sommario}{Impossibile numerare la bibliografia}{}}

\setstretch{1.15}

\pagestyle{fancy}
\fancyhf{}
\renewcommand{\headrulewidth}{0.4pt}
\fancyhead[L]{\raisebox{-0.28\height}{\includegraphics[height=1.02cm]{ucbm-logo.png}}}
\fancyhead[R]{\parbox[c]{11.0cm}{\raggedleft\footnotesize\itshape
  Sommario del Lavoro di Tesi, Laurea Magistrale in Ingegneria dei
  Sistemi Intelligenti, Universit\`a Campus Bio-Medico di Roma}}
\fancyfoot[C]{\footnotesize\thepage}

\setlength{\abovedisplayskip}{4pt plus 2pt minus 2pt}
\setlength{\belowdisplayskip}{4pt plus 2pt minus 2pt}
\setlength{\abovedisplayshortskip}{2pt plus 1pt}
\setlength{\belowdisplayshortskip}{3pt plus 1pt}
\setlength{\parindent}{0pt}
\setlength{\parskip}{3pt plus 0.5pt minus 0.5pt}
\setlist{nosep,leftmargin=*,topsep=2pt,itemsep=1pt}
\makeatletter
\renewcommand\section{\@startsection{section}{1}{\z@}%
  {-1.4ex \@plus -.4ex \@minus -.2ex}{0.4ex \@plus .1ex}%
  {\normalfont\large\bfseries}}
\renewcommand\subsection{\@startsection{subsection}{2}{\z@}%
  {-1.0ex \@plus -.3ex \@minus -.2ex}{0.1ex \@plus .1ex}%
  {\normalfont\normalsize\bfseries}}
\makeatother

\setlength{\abovecaptionskip}{4pt}
\setlength{\belowcaptionskip}{2pt}
\setlength{\textfloatsep}{8pt plus 2pt minus 2pt}
\setlength{\intextsep}{8pt plus 2pt minus 2pt}
\renewcommand{\arraystretch}{0.98}

\newcommand{\titoloIT}{Shor's Algorithm under Noise: Machine-Learning Ablation,
  Quantum Error Correction, and Approximate QFT}
\newcommand{\laureando}{Claudio Dragotta}
\newcommand{\relatore}{Prof. Paolo Soda}
\newcommand{\correlatore}{Ing. Floriano Caprio}
\newcommand{\sessione}{Ottobre --- A.A. 2025/2026}
\newcommand{\intestazioneIT}{Sommario del Lavoro di Tesi, Laurea Magistrale in Ingegneria dei
  Sistemi Intelligenti, Universit\`a Campus Bio-Medico di Roma}
\newcommand{\intestazioneEN}{Thesis Summary, Master's Degree in Intelligent Systems Engineering,
  Universit\`a Campus Bio-Medico di Roma}

\newcommand{\bloccoTitolo}[6]{%
  \begin{center}
    {\large\bfseries #1}%
    \if\relax\detokenize{#2}\relax\else
      \\[3pt]{\small\itshape #2}%
    \fi
  \end{center}
  \vspace{-4pt}
  {\small\setlength{\tabcolsep}{0pt}%
  \begin{tabular}{p{0.46\textwidth}p{0.53\textwidth}}
    \textbf{#3:} \laureando & \textbf{#4:} \relatore\\[1pt]
    \textbf{#5:} \correlatore & \textbf{#6:} \sessione
  \end{tabular}}
  \vspace{3pt}
  \hrule
  \vspace{-2pt}
}

\begin{document}

\bloccoTitolo{\titoloIT}{}{Laureando}{Relatore}{Correlatore}{Sessione di Laurea}

\section{Introduzione e obiettivi}

L'algoritmo di Shor fattorizza numeri interi in tempo polinomiale rispetto alla lunghezza
binaria dell'input, con implicazioni per la sicurezza dei sistemi crittografici basati sulla
fattorizzazione \cite{shor1994,nist2024pqc,gidney2021}. Errori nelle operazioni, decoerenza e
lettura imperfetta possono però compromettere l'informazione necessaria alla ricerca
dell'ordine moltiplicativo \cite{preskill2018}.

La tesi ne studia la sensibilità al rumore e valuta tre strategie: selezione degli esiti
mediante post-processing, isolando il contributo del machine learning; correzione quantistica
degli errori (QEC) con surface code e codici qLDPC; riduzione delle porte mediante trasformata
di Fourier quantistica approssimata (AQFT). Ogni strategia è confrontata con riferimenti
espliciti e metriche coerenti con il compito. Le campagne simulative sono separate: la QEC è
valutata su circuiti di memoria, non integrata nell'intero algoritmo di Shor.

\section{Materiali e metodi}

\subsection{Shor, validazione e definizione del successo}\label{sec:implementazione}

Scelta una base \(a\) coprima con il numero \(N\) da fattorizzare, l'ordine \(r\), ossia il
periodo di \(a^{x} \bmod N\), è il minimo intero positivo per cui \(a^{r} \equiv 1 \pmod{N}\).
La parte quantistica lo stima mediante stima di fase e trasformata di Fourier quantistica
inversa (QFT inversa); quella classica ricava e verifica i fattori mediante approssimazione
razionale e massimo comune divisore.

L'istanza principale, \(N = 15\) e \(a = 7\), usa otto qubit di controllo e quattro di lavoro.
Con \(r = 4\), gli esiti ideali sono 0, 64, 128 e 192. Ogni esecuzione con misura
(\emph{shot}) produce un esito; molti shot formano un istogramma. Il successo è il recupero di
fattori non banali, non dell'ordine esatto: tre picchi ideali vi conducono, con probabilità
del 75\%. Il post-processing accetta tuttavia 63 dei 256 esiti: una distribuzione uniforme
darebbe quindi il 24,61\% di successo. Questo riferimento dipende dall'istanza e dalla
verifica classica, non misura la fedeltà dello stato quantistico \cite{smolin2013}.

La validazione controlla separatamente operatori modulari, ripristino degli ausiliari e
distribuzioni teoriche. Sono validate idealmente anche \(N = 21\), \(a = 2\) e \(N = 35\),
\(a = 6\), con aritmetica reversibile di Beauregard \cite{beauregard2002}; solo \(N = 21\) entra
nelle prove rumorose esplorative sull'AQFT, per limiti delle risorse di simulazione. Si
utilizzano Qiskit/Aer, scikit-learn, Stim, PyMatching e ldpc, registrando versioni,
parametri, semi casuali e identificativi circuitali.

\subsection{Sensibilità agli errori di porta}

Il circuito principale è compilato nelle porte RZ, SX, X e CX: profondità 412, con 224 CX e
70 SX. Dopo ciascuna porta, escluse le RZ virtuali, si introducono errori di Pauli
indipendenti, variando la probabilità di errore non identità \(p_g\) su 13 valori tra 0 e 0,5.
Il modello è fenomenologico e uniforme, non calibrato su un dispositivo reale. La probabilità
per porta \(p_g\) non è direttamente confrontabile con il fallimento logico \(p_L\) dei codici
QEC: Shor non è compilato in forma fault-tolerant.

\subsection{Post-processing e ablazione del machine learning}

Tre criteri selezionano gli esiti da verificare sugli stessi istogrammi da 1.024 shot, con
estrattore dei fattori e gestione dei pareggi comuni. Ogni istogramma costituisce
un'iterazione, ripetuta in caso di insuccesso. TOP-1 verifica il candidato più frequente;
TOP-4 ne verifica fino a quattro, ordinati per frequenza; il Metodo~2 applica TOP-4 solo se un
classificatore predice il successo del candidato dominante. Il confronto con TOP-4 senza
filtro costituisce l'ablazione: separa il contributo del modello appreso dal beneficio di
provare più candidati.

Gli scenari illustrativi UC1, di riferimento, e UC2, di stress, includono depolarizzazione,
rilassamento termico ed errori di lettura. Per scenario, 2.000 istogrammi, con parametri
variati entro \(\pm50\%\) del nominale, sono divisi 60\%/20\%/20\% tra addestramento,
selezione e verifica. Random forest, support vector machine (SVM) e perceptron multistrato
sono confrontati sul punteggio F1; il modello scelto è valutato su 400 istogrammi
indipendenti. La metrica primaria è il numero di iterazioni per ottenere i fattori, con 30
repliche appaiate e un massimo di 50 iterazioni; l'accuratezza è secondaria e i tempi
dipendono dalle risorse disponibili.

\subsection{Correzione quantistica: surface code e qLDPC}

La QEC distribuisce l'informazione di un qubit logico su più qubit fisici. Le sindromi
segnalano gli errori senza misurare direttamente lo stato logico e guidano il decoder nella
correzione; la metrica è la probabilità di fallimento logico \(p_L\). Dopo verifiche
preliminari con codice a ripetizione e codice di Steane, si studiano memorie surface code
\cite{steane1996,fowler2012,googleQECBelowThreshold2025} a distanza \(d\) crescente, parametro
legato alla capacità correttiva, con errori di operazione, preparazione e misura e, in alcune
configurazioni, crosstalk tra qubit adiacenti.

Il decoder di riferimento, \emph{minimum-weight perfect matching} (MWPM), associa a coppie le
sindromi attivate. È confrontato con una rete neurale autonoma e con un decoder ibrido, in cui
una rete decide quando modificare la scelta di MWPM. La soglia d'intervento è selezionata sul
validation set e fissata sul test set \cite{alphaqubit2024}. Poiché i circuiti di correzione
contengono solo porte di Clifford e misure, sono simulabili in modo efficiente su un computer
classico: Stim ne campiona le sindromi per milioni di esecuzioni e i decoder, compresa la rete
neurale, sono addestrati e valutati su questi dati.

Un confronto esplorativo considera il Gross code di IBM, un codice quantum Low-Density
Parity-Check (qLDPC) con 12 qubit logici e 144 qubit di dato \cite{bravyi2024}. Si usa il
modello \emph{code-capacity}: errori sui soli qubit di dato e sindromi perfette. BP+OSD
(\emph{belief propagation with ordered statistics decoding}) è impiegato sul Gross code e come
ulteriore riferimento sul surface code \cite{roffe2020}. Per il Gross code si confrontano più
configurazioni, inclusi aggiornamenti dei messaggi paralleli o seriali.

\subsection{Riduzione del costo circuitale mediante AQFT}

L'AQFT elimina le rotazioni controllate di piccolo angolo, riducendo porte e profondità a
scapito della precisione di fase \cite{barencoAQFT1996}. Si valuta se il rumore evitato
compensi questa perdita. È applicata alla QFT inversa finale di Shor per \(N = 15\) e alle QFT
interne all'aritmetica modulare per \(N = 21\); benchmark separati di stima di fase, su
registri da sei a dieci qubit, ne isolano l'effetto.

Il rumore deriva dalle infedeltà delle porte e dagli errori di lettura della calibrazione
offline FakeSherbrooke, senza contributi aggiuntivi \(T_1/T_2\); layout e semi casuali restano
invariati. Il grado di approssimazione, cioè quante rotazioni eliminare, è scelto su dati di
selezione e valutato su dati indipendenti. Gli 8.192 shot per configurazione del pilota
\(N = 15\) e i 256 per \(N = 21\) sono divisi equamente tra selezione e verifica.

\section{Risultati}

\subsection{Validazione e sensibilità di Shor al rumore}

Le verifiche ideali confermano la correttezza dei circuiti; la distanza di variazione totale
dalla distribuzione teorica è \(1{,}38 \times 10^{-9}\) per \(N = 21\) e
\(2{,}50 \times 10^{-13}\) per \(N = 35\). Per ogni \(p_g\) si eseguono 20 repliche da 4.096
shot. Il successo senza rumore è 74,89\% (intervallo di confidenza al 95\%
{[}74,59\%; 75,18\%{]}), coerente con il 75\% atteso. Scende a 72,5\%, 56,5\%, 45,2\% e
29,5\% per \(p_g\) rispettivamente pari a 0,1\%, 1\%, 2\% e 5\%. Per \(p_g \geq 10\%\) resta
intorno al 25\%, fino al 24,59\% per \(p_g = 0{,}5\), prossimo al riferimento uniforme del
24,61\%. Con 294 porte esposte, \(p_g = 1\%\) implica circa tre errori medi per esecuzione.

\subsection{Post-processing e machine learning}

Le SVM selezionate raggiungono, in UC1 e UC2, F1 pari a 0,919 e 0,923 e area sotto la curva
ROC pari a 0,965 e 0,958. Questa capacità predittiva non riduce però le iterazioni. TOP-4
ottiene i fattori alla prima iterazione nelle repliche osservate di entrambi gli scenari:
media 1,00, contro 1,37 e 1,50 di TOP-1, con differenze significative nei test appaiati. Il
Metodo~2 richiede invece 1,50 e 1,37 iterazioni: non migliora significativamente TOP-1 ed è
peggiore di TOP-4. Il beneficio deriva quindi dai candidati aggiuntivi, non dal
classificatore, nelle condizioni studiate.

\subsection{Correzione quantistica degli errori}\label{sec:risultati-qec}

Indicando con \(p\) la probabilità di errore fisico, il codice a ripetizione riproduce, entro
l'incertezza Monte Carlo, \(p_L = 3p^2 - 2p^3\); Steane mostra una dipendenza circa
quadratica dall'errore fisico, con pendenza logaritmica 1,94 e pseudo-soglia circa 0,08. Nel
surface code, aumentare la distanza riduce l'errore logico sotto soglia e lo aumenta sopra
soglia. Il fit stima soglie circa 0,86\% in base \(Z\) e 0,81\% in base \(X\), nel modello
considerato.

Con rumore nominale la rete autonoma non supera significativamente MWPM. Con crosstalk, per
\(d = 3\), \(p = 0{,}003\) e crosstalk 0,01, il decoder ibrido riduce \(p_L\) da 0,05829 a
0,04851: \(-16{,}8\%\) relativo, significativo al test appaiato di McNemar. Il beneficio si
riduce a \(d = 5\) e non è significativo a \(d = 7\); nelle configurazioni a distanza 5 BP+OSD
supera l'ibrido. Le prestazioni dipendono anche dai cicli di correzione e dalla dimensione
dell'ingresso della rete.

Nel modello code-capacity, il Gross code protegge 12 qubit logici con 144 qubit di dato meglio
di 12 patch indipendenti di surface code a distanza 11 (1.452 qubit di dato) a tutti i livelli
di rumore studiati, e di patch a distanza 13 (2.028) per \(p \geq 0{,}75\%\). A \(p = 1\%\) il
fallimento di blocco è \(1{,}0 \times 10^{-6}\), contro \(1{,}6 \times 10^{-6}\) delle patch a
distanza 13. Con aggiornamento seriale, il Gross code corregge tutti gli errori fino a cinque
qubit, verificati su circa \(5 \times 10^{8}\) configurazioni; a \(p = 1\%\) l'errore logico è
circa cento volte inferiore rispetto alla configurazione parallela iniziale, che falliva già
con tre errori. Il vantaggio dipende dunque dal decoder, nel modello studiato.

\subsection{Effetti dell'AQFT}

Per \(N = 15\), sui dati di verifica l'AQFT selezionata aumenta il successo per shot da
47,92\% a 62,79\%: +14,87 punti percentuali (+31\% relativo), con intervallo di confidenza
Newcombe al 95\% {[}12,73; 16,99{]}. Le porte a due qubit ECR scendono da 362 a 279
(\(-22{,}9\%\)) e la profondità da 1.351 a 1.141 (\(-15{,}5\%\)). Il beneficio riguarda
direttamente la fattorizzazione.

Per \(N = 21\), le ECR si riducono da 49.661 a 23.178 (\(-53{,}3\%\)), ma il successo ideale
scende da 46,67\% a 40,13\% (\(-6{,}54\) punti percentuali). Nei tre confronti rumorosi, con
128 shot di verifica per variante, gli intervalli delle differenze comprendono lo zero: un
beneficio non è dimostrato. Nei nove benchmark di stima di fase tutte le differenze favoriscono
l'AQFT, fino a 12,55 punti percentuali ({[}9,77; 15,30{]}). I confronti restano esplorativi:
il grado di approssimazione migliore tra quelli provati, due o tre, non è generalizzabile a
registri arbitrari.

\section{Discussione e conclusioni}

Le strategie agiscono su livelli distinti. Il post-processing sfrutta l'informazione residua
senza modificare lo stato quantistico: nell'ablazione, buone metriche del classificatore non
implicano un minor costo di fattorizzazione. Alle alte intensità di rumore, il successo
prossimo al riferimento uniforme non dimostra che l'informazione quantistica sia preservata.

La QEC protegge l'informazione e, negli esperimenti di memoria, la decodifica appresa può
essere utile quando le sindromi contengono correlazioni non rappresentate dal decoder di
riferimento, come il crosstalk. Il confronto qLDPC mostra che struttura del codice e decoder
determinano insieme il rapporto tra protezione e numero di qubit di dato; non costituisce una
misura del miglioramento di Shor.

L'AQFT riduce le opportunità di errore. Per \(N = 15\), l'ordine \(r = 4\) genera fasi
rappresentabili esattamente con due bit; per \(N = 21\), \(r = 6\) richiede fasi non
esattamente rappresentabili e un'aritmetica più complessa. Il beneficio non è automatico,
l'approssimazione va scelta per istanza e rumore e verificata su dati indipendenti.

Le conclusioni restano simulative: post-processing su \(N = 15\), QEC su memorie separate,
Gross code in code-capacity e AQFT su poche istanze con una calibrazione offline. Non si
realizza Shor fault-tolerant né si stimano risorse per la scala crittografica.

Per gli sviluppi futuri si propone pertanto di privilegiare la riduzione controllata del costo
circuitale, preservando l'informazione necessaria alla ricerca dell'ordine. Oltre all'AQFT
\cite{barencoAQFT1996}, andrebbero investigate ottimizzazioni dell'aritmetica modulare e della
sintesi delle operazioni \cite{gidney2021}. La qualità dovrà essere verificata mediante
accuratezza della stima di fase e successo della fattorizzazione, non dal solo numero di porte.
Non si dimostra qui una modifica della complessità asintotica di Shor.

Questa direzione va integrata con una futura implementazione fault-tolerant, basata su qubit
logici e operazioni protette che limitino la propagazione degli errori. Ottimizzazione e QEC
sono complementari: la prima mira a contenere le risorse, la seconda a proteggere l'esecuzione
dal rumore.

\par
\vspace{-4pt}
\section{Riferimenti bibliografici essenziali}
\begingroup
\setstretch{1}
\fontsize{6.4}{7.0}\selectfont
\setlength{\columnsep}{10pt}
\begin{multicols}{2}
\raggedcolumns
\raggedright %
\renewcommand{\section}[1]{}%
\begin{thebibliography}{14}
\setlength{\itemsep}{-1pt}
\setlength{\parsep}{0pt}
\bibitem{shor1994} P. W. Shor, ``Algorithms for Quantum Computation: Discrete Logarithms and
Factoring'', \emph{Proc. 35th Annual Symposium on Foundations of Computer Science (FOCS)},
pp.~124--134, 1994.

\bibitem{nist2024pqc} National Institute of Standards and Technology,
\emph{FIPS 203, FIPS 204 e FIPS 205}, 2024.

\bibitem{gidney2021} C. Gidney e M. Eker\aa, ``How to Factor 2048 Bit RSA Integers in
8 Hours Using 20 Million Noisy Qubits'', \emph{Quantum}, 5, 433, 2021.

\bibitem{preskill2018} J. Preskill, ``Quantum Computing in the NISQ Era and Beyond'', \emph{Quantum},
2, 79, 2018.

\bibitem{smolin2013} J. A. Smolin, G. Smith e A. Vargo, ``Oversimplifying Quantum
Factoring'', \emph{Nature}, 499, pp.~163--165, 2013.

\bibitem{beauregard2002} S. Beauregard, ``Circuit for Shor's Algorithm Using \(2n+3\) Qubits'',
\emph{Quantum Inf. Comput.}, 3(2), pp.~175--185, 2003.

\bibitem{steane1996} A. M. Steane, ``Error Correcting Codes in Quantum Theory'',
\emph{Phys. Rev. Lett.}, 77(5), pp.~793--797, 1996.

\bibitem{fowler2012} A. G. Fowler et al., ``Surface Codes: Towards Practical Large-Scale
Quantum Computation'', \emph{Phys. Rev. A}, 86, 032324, 2012.

\bibitem{googleQECBelowThreshold2025} Google Quantum AI and Collaborators, ``Quantum
Error Correction below the Surface Code Threshold'', \emph{Nature}, 638,
pp.~920--926, 2025.

\bibitem{alphaqubit2024} J. Bausch et al., ``Learning High-Accuracy Error Decoding for
Quantum Processors'', \emph{Nature}, 635, pp.~834--840, 2024.

\bibitem{bravyi2024} S. Bravyi et al., ``High-Threshold and Low-Overhead Fault-Tolerant
Quantum Memory'', \emph{Nature}, 627, pp.~778--782, 2024.

\bibitem{roffe2020} J. Roffe et al., ``Decoding Across the Quantum Low-Density Parity-Check Code
Landscape'', \emph{Phys. Rev. Research}, 2, 043423, 2020.

\bibitem{barencoAQFT1996} A. Barenco, A. Ekert, K.-A. Suominen e P. Törmä,
``Approximate Quantum Fourier Transform and Decoherence'', \emph{Phys. Rev. A},
54, pp.~139--146, 1996.
\end{thebibliography}
\end{multicols}
\endgroup


\clearpage
\selectlanguage{english}
\setcounter{section}{0}
\renewcommand{\theHsection}{en.\arabic{section}}
\renewcommand{\theHsubsection}{en.\arabic{section}.\arabic{subsection}}
\fancyhead[R]{\parbox[c]{11.0cm}{\raggedleft\footnotesize\itshape
  \intestazioneEN}}
\renewcommand{\sessione}{October --- A.Y. 2025/2026}
\bloccoTitolo{\titoloIT}{}{Candidate}{Supervisor}{Co-supervisor}{Session}

\section{Introduction and objectives}

Shor's algorithm factors integers in polynomial time with respect to the binary length of
the input, with implications for the security of cryptographic systems based on
factorization \cite{en-shor1994,en-nist2024pqc,en-gidney2021}. Errors in operations,
decoherence, and imperfect readout, however, can compromise the information needed to find
the multiplicative order \cite{en-preskill2018}.

The thesis studies its sensitivity to noise and evaluates three strategies: outcome
selection through post-processing, isolating the contribution of machine learning; quantum
error correction (QEC) with surface codes and qLDPC codes; gate reduction through the
approximate quantum Fourier transform (AQFT). Each strategy is compared with explicit
references and with metrics consistent with the task. The simulation campaigns are
separate: QEC is evaluated on memory circuits, not within the full Shor algorithm.

\section{Materials and methods}

\subsection{Shor's algorithm, validation, and definition of success}\label{sec:implementation-en}

Given a base \(a\) coprime to the number \(N\) to be factored, the order \(r\), i.e.\ the
period of \(a^{x} \bmod N\), is the smallest positive integer such that
\(a^{r} \equiv 1 \pmod{N}\). The quantum part estimates it through phase estimation and the
inverse quantum Fourier transform (inverse QFT); the classical part derives and verifies
the factors through rational approximation and the greatest common divisor.

The main instance, \(N = 15\) and \(a = 7\), uses eight control qubits and four work
qubits. With \(r = 4\), the ideal outcomes are 0, 64, 128, and 192. Each measured execution
(\emph{shot}) yields one outcome; many shots form a histogram. Success is the recovery of
non-trivial factors, not of the exact order: three ideal peaks lead to it, with a
probability of 75\%. Post-processing, however, accepts 63 of the 256 outcomes, so a uniform
distribution would still give 24.61\% success. This reference depends on the instance and
on the classical verification, and does not measure the fidelity of the quantum state
\cite{en-smolin2013}.

Validation separately checks modular operators, ancilla restoration, and theoretical
distributions. \(N = 21\), \(a = 2\) and \(N = 35\), \(a = 6\) are also validated ideally,
with Beauregard's reversible arithmetic \cite{en-beauregard2002}; only \(N = 21\) enters the
exploratory noisy AQFT tests, owing to limited simulation resources. Qiskit/Aer,
scikit-learn, Stim, PyMatching, and ldpc are used, recording versions, parameters, seeds,
and circuit identifiers.

\subsection{Sensitivity to gate errors}

The main circuit is compiled into RZ, SX, X, and CX gates: depth 412, with 224 CX and 70
SX. After each gate, excluding the virtual RZ gates, independent Pauli errors are
introduced, varying the non-identity error probability \(p_g\) over 13 values between 0 and
0.5. The model is phenomenological and uniform, not calibrated on a real device. The
per-gate probability \(p_g\) is not directly comparable with the logical failure \(p_L\) of
QEC codes: Shor's algorithm is not compiled in fault-tolerant form.

\subsection{Post-processing and machine-learning ablation}

Three criteria select the outcomes to verify on the same 1,024-shot histograms, with a
common factor extractor and tie handling. Each histogram constitutes one iteration,
repeated in case of failure. TOP-1 verifies the most frequent candidate; TOP-4 verifies up
to four, ranked by frequency; Method~2 applies TOP-4 only if a classifier predicts success
for the dominant candidate. The comparison with unfiltered TOP-4 constitutes the ablation:
it separates the contribution of the learned model from the benefit of trying more
candidates.

The illustrative scenarios UC1, the reference, and UC2, the stress case, include
depolarization, thermal relaxation, and readout errors. For each scenario, 2,000
histograms, with parameters varied within \(\pm50\%\) of nominal, are split 60\%/20\%/20\%
into training, selection, and test. Random forests, support vector machines (SVMs), and
multilayer perceptrons are compared on the F1 score; the chosen model is evaluated on 400
independent histograms. The primary metric is the number of iterations needed to obtain the
factors, with 30 paired replicates and at most 50 iterations; accuracy is secondary and run
times depend on the available resources.

\subsection{Quantum error correction: surface code and qLDPC}

QEC distributes the information of one logical qubit over several physical qubits.
Syndromes flag errors without directly measuring the logical state and guide the decoder in
the correction; the metric is the logical failure probability \(p_L\). After preliminary
checks with the repetition code and the Steane code, surface-code memories
\cite{en-steane1996,en-fowler2012,en-googleQECBelowThreshold2025} are studied at increasing
distance \(d\), a parameter linked to the correction capability, with operation,
preparation, and measurement errors and, in some configurations, crosstalk between adjacent
qubits.

The reference decoder, \emph{minimum-weight perfect matching} (MWPM), pairs the triggered
syndromes. It is compared with a stand-alone neural network and with a hybrid decoder, in
which a network decides when to override the MWPM choice. The intervention threshold is
selected on the validation set and fixed on the test set \cite{en-alphaqubit2024}. Since
correction circuits contain only Clifford gates and measurements, they can be simulated
efficiently on a classical computer: Stim samples their syndromes over millions of runs,
and the decoders, including the neural network, are trained and evaluated on these data.

An exploratory comparison considers IBM's Gross code, a quantum low-density parity-check
(qLDPC) code with 12 logical qubits and 144 data qubits \cite{en-bravyi2024}. The
\emph{code-capacity} model is used: errors on data qubits only and perfect syndromes.
BP+OSD (\emph{belief propagation with ordered statistics decoding}) is used on the Gross
code and as a further reference on the surface code \cite{en-roffe2020}. For the Gross
code, several configurations are compared, including parallel and serial message updates.

\subsection{Reducing circuit cost with the AQFT}

The AQFT removes small-angle controlled rotations, reducing gates and depth at the expense
of phase precision \cite{en-barencoAQFT1996}. The question is whether the noise avoided
compensates for this loss. It is applied to Shor's final inverse QFT for \(N = 15\) and to
the QFTs inside the modular arithmetic for \(N = 21\); separate phase-estimation
benchmarks, on registers of six to ten qubits, isolate its effect.

The noise is derived from the gate infidelities and readout errors of the offline
FakeSherbrooke calibration, without additional \(T_1/T_2\) contributions; layouts and
random seeds are kept unchanged. The approximation degree, i.e.\ how many rotations to
remove, is chosen on selection data and evaluated on independent data. The 8,192 shots per
configuration of the \(N = 15\) pilot and the 256 for \(N = 21\) are split equally between
selection and test.

\section{Results}

\subsection{Validation and sensitivity of Shor's algorithm to noise}

Ideal checks confirm that the circuits are correct; the total variation distance from the
theoretical distribution is \(1.38 \times 10^{-9}\) for \(N = 21\) and
\(2.50 \times 10^{-13}\) for \(N = 35\). For each \(p_g\), 20 replicates of 4,096 shots are
run. Noise-free success is 74.89\% (95\% confidence interval {[}74.59\%; 75.18\%{]}),
consistent with the expected 75\%. It falls to 72.5\%, 56.5\%, 45.2\%, and 29.5\% for
\(p_g\) equal to 0.1\%, 1\%, 2\%, and 5\%, respectively. For \(p_g \geq 10\%\) it stays
around 25\%, down to 24.59\% at \(p_g = 0.5\), close to the uniform reference of 24.61\%.
With 294 exposed gates, \(p_g = 1\%\) implies about three errors on average per run.

\subsection{Post-processing and machine learning}

The selected SVMs reach, in UC1 and UC2, F1 scores of 0.919 and 0.923 and areas under the
ROC curve of 0.965 and 0.958. This predictive ability, however, does not reduce the
iterations. TOP-4 obtains the factors at the first iteration in the observed replicates of
both scenarios: mean 1.00, versus 1.37 and 1.50 for TOP-1, with significant differences in
paired tests. Method~2 instead requires 1.50 and 1.37 iterations: it does not
significantly improve on TOP-1 and is worse than TOP-4. The benefit thus comes from the
additional candidates, not the classifier, under the conditions studied.

\subsection{Quantum error correction}\label{sec:results-qec-en}

Denoting by \(p\) the physical error probability, the repetition code reproduces, within
Monte Carlo uncertainty, \(p_L = 3p^2 - 2p^3\); the Steane code shows an approximately
quadratic dependence on the physical error, with a logarithmic slope of 1.94 and a
pseudo-threshold of about 0.08. In the surface code, increasing the distance reduces the
logical error below threshold and increases it above threshold. The fit estimates
thresholds of about 0.86\% in the \(Z\) basis and 0.81\% in the \(X\) basis, in this model.

Under nominal noise, the stand-alone network does not significantly outperform MWPM. With
crosstalk, for \(d = 3\), \(p = 0.003\), and crosstalk 0.01, the hybrid decoder reduces
\(p_L\) from 0.05829 to 0.04851: \(-16.8\%\) relative, significant by the paired McNemar
test. The benefit shrinks at \(d = 5\) and is not significant at \(d = 7\); in the
distance-5 configurations BP+OSD outperforms the hybrid. Performance also depends on the
number of correction rounds and on the size of the network input.

In the code-capacity model, the Gross code protects 12 logical qubits with 144 data qubits
better than 12 independent surface-code patches at distance 11 (1,452 data qubits) at all
noise levels studied, and better than distance-13 patches (2,028) for
\(p \geq 0.75\%\). At \(p = 1\%\) the block failure is \(1.0 \times 10^{-6}\), versus
\(1.6 \times 10^{-6}\) for the distance-13 patches. With serial updates, the Gross code
corrects all errors of up to five qubits, verified over about \(5 \times 10^{8}\)
configurations; at \(p = 1\%\) the logical error is about a hundred times lower than with
the initial parallel configuration, which already failed with three errors. The advantage
therefore depends on the decoder, in this model.

\subsection{Effects of the AQFT}

For \(N = 15\), on the test data the selected AQFT raises per-shot success from 47.92\% to
62.79\%: +14.87 percentage points (+31\% relative), with a 95\% Newcombe confidence
interval of {[}12.73; 16.99{]}. Two-qubit ECR gates decrease from 362 to 279
(\(-22.9\%\)) and depth from 1,351 to 1,141 (\(-15.5\%\)). The benefit concerns
factorization directly.

For \(N = 21\), ECR gates drop from 49,661 to 23,178 (\(-53.3\%\)), but ideal success falls
from 46.67\% to 40.13\% (\(-6.54\) percentage points). In the three noisy comparisons, with
128 test shots per variant, the intervals of the differences include zero: a benefit is
not demonstrated. In the nine phase-estimation benchmarks all differences favor the AQFT,
up to 12.55 percentage points ({[}9.77; 15.30{]}). The comparisons remain exploratory: the
best approximation degree among those tested, two or three, cannot be generalized to
arbitrary registers.

\section{Discussion and conclusions}

The strategies act at distinct levels. Post-processing exploits the residual information
without modifying the quantum state: in the ablation, good classifier metrics do not imply
a lower factorization cost. At high noise intensities, success close to the uniform
reference does not show that quantum information is preserved.

QEC protects information; in memory experiments, learned decoding can help when syndromes
contain correlations not represented by the reference decoder, such as crosstalk. The
qLDPC comparison shows that code structure and decoder jointly set the balance between
protection and data qubits, without measuring gains for Shor's algorithm.

The AQFT reduces opportunities for error. For \(N = 15\), the order \(r = 4\) generates
phases exactly representable with two bits; for \(N = 21\), \(r = 6\) requires phases not
exactly representable and more complex arithmetic. The benefit is not automatic; the
approximation must be chosen per instance and noise, and verified on independent data.

The conclusions remain simulation-based: post-processing on \(N = 15\), QEC on separate
memories, the Gross code in code-capacity, and the AQFT on few instances with one offline
calibration. No fault-tolerant execution of Shor's algorithm is realized, and no resources
are estimated for the cryptographic scale.

Future work should therefore prioritize the controlled reduction of circuit cost,
preserving the information needed for order finding. Besides the AQFT
\cite{en-barencoAQFT1996}, optimizations of modular arithmetic and operation synthesis
should be studied \cite{en-gidney2021}. Quality must be verified through phase-estimation
accuracy and factorization success, not gate count alone. No change in Shor's asymptotic
complexity is shown here.

This direction should accompany future fault-tolerant execution, with logical qubits and
protected operations limiting error propagation. Optimization and QEC are complementary:
the former contains resources, the latter protects execution from noise.

\par
\vspace{-4pt}
\section{Essential references}
\begingroup
\setstretch{1}
\fontsize{6.4}{7.0}\selectfont
\setlength{\columnsep}{10pt}
\begin{multicols}{2}
\raggedcolumns
\raggedright
\renewcommand{\section}[1]{}%
\begin{thebibliography}{14}
\setlength{\itemsep}{-1pt}
\setlength{\parsep}{0pt}
\bibitem{en-shor1994} P. W. Shor, ``Algorithms for Quantum Computation: Discrete Logarithms and
Factoring'', \emph{Proc. 35th Annual Symposium on Foundations of Computer Science (FOCS)},
pp.~124--134, 1994.

\bibitem{en-nist2024pqc} National Institute of Standards and Technology,
\emph{FIPS 203, FIPS 204, and FIPS 205}, 2024.

\bibitem{en-gidney2021} C. Gidney and M. Eker\aa, ``How to Factor 2048 Bit RSA Integers
in 8 Hours Using 20 Million Noisy Qubits'', \emph{Quantum}, 5, 433, 2021.

\bibitem{en-preskill2018} J. Preskill, ``Quantum Computing in the NISQ Era and Beyond'',
\emph{Quantum}, 2, 79, 2018.

\bibitem{en-smolin2013} J. A. Smolin, G. Smith, and A. Vargo, ``Oversimplifying Quantum
Factoring'', \emph{Nature}, 499, pp.~163--165, 2013.

\bibitem{en-beauregard2002} S. Beauregard, ``Circuit for Shor's Algorithm Using \(2n+3\) Qubits'',
\emph{Quantum Inf. Comput.}, 3(2), pp.~175--185, 2003.

\bibitem{en-steane1996} A. M. Steane, ``Error Correcting Codes in Quantum Theory'',
\emph{Phys. Rev. Lett.}, 77(5), pp.~793--797, 1996.

\bibitem{en-fowler2012} A. G. Fowler et al., ``Surface Codes: Towards Practical
Large-Scale Quantum Computation'', \emph{Phys. Rev. A}, 86, 032324, 2012.

\bibitem{en-googleQECBelowThreshold2025} Google Quantum AI and Collaborators,
``Quantum Error Correction below the Surface Code Threshold'', \emph{Nature}, 638,
pp.~920--926, 2025.

\bibitem{en-alphaqubit2024} J. Bausch et al., ``Learning High-Accuracy Error Decoding
for Quantum Processors'', \emph{Nature}, 635, pp.~834--840, 2024.

\bibitem{en-bravyi2024} S. Bravyi et al., ``High-Threshold and Low-Overhead Fault-Tolerant
Quantum Memory'', \emph{Nature}, 627, pp.~778--782, 2024.

\bibitem{en-roffe2020} J. Roffe et al., ``Decoding Across the Quantum Low-Density Parity-Check Code
Landscape'', \emph{Phys. Rev. Research}, 2, 043423, 2020.

\bibitem{en-barencoAQFT1996} A. Barenco, A. Ekert, K.-A. Suominen, and P. Törmä,
``Approximate Quantum Fourier Transform and Decoherence'', \emph{Phys. Rev. A},
54, pp.~139--146, 1996.
\end{thebibliography}
\end{multicols}
\endgroup


\end{document}
